% articulo-08-indice.tex - Índice de contenido.
% (tocloft)
%
% Diseño:
%   · Todo en cuerpo normal (12 pt).
%   · Nivel 1 (section): MAYÚSCULAS, número al margen + folio.
%   · Nivel 2 (subsection): redonda, sin número, con folio.
%   · Nivel 3 (subsubsection): cursiva, sin número, sangría 0.5\ParIndent, con folio.
%   · \tocblockskip separa bloques (preliminares / cuerpo / finales).
%
% Uso en el .md:
%   Separar bloques:   ```{=latex}\n\addtocontents{toc}{\protect\tocblockskip}\n```
%   Entrada sin folio: ```{=latex}\n\addcontentsline{toc}{part}{Título}\n```
%
% Variables YAML:
%   numbersections: true   Activa numeración de secciones (por defecto desactivada).
%   toc-depth: 3           Profundidad del índice (por defecto 3).
%   titulos-sans: true     El índice completo va en sans-serif, coherente con la
%                          opción homónima de 13-secciones.tex y 07-encabezado.tex.
%                          Afecta a los tres niveles de entrada y a los folios.
%   (sin declarar)         Índice en la familia principal del documento (por defecto).

\usepackage[titles]{tocloft}

% Numeración de secciones: desactivada por defecto.

$if(numbersections)$
\setcounter{secnumdepth}{5}
$else$
\setcounter{secnumdepth}{-\maxdimen}
$endif$

% Profundidad del índice: defecto 3 (section, subsection, subsubsection).

\setcounter{tocdepth}{$if(toc-depth)$$toc-depth$$else$3$endif$}

\newlength{\tocsubindent}
\setlength{\tocsubindent}{0.5\ParIndent}   % sangría nivel 3

\newcommand{\tocblockskip}{\addvspace{\baselineskip}}

\setlength{\cftbeforesecskip}{0.4\baselineskip}
\setlength{\cftbeforesubsecskip}{0pt}
\setlength{\cftbeforesubsubsecskip}{0pt}

\makeatletter

% ── \toc@line{indent}{fontcmd}{number}{title}{page} ──────────────────────────
%
% Se controla el margen con \leftskip (no con \noindent+\hspace+\parbox).
% \noindent suprime solo \parindent, pero \leftskip puede estar heredado de
% algún entorno padre y seguir desplazando el texto. Fijando explícitamente
% \leftskip=#1, \parindent=0pt y \parfillskip=0pt plus 1fil (raggedright)
% el resultado es independiente del contexto en que se invoque.
%
% Con titulos-sans: true, \toc@font conmuta a \sffamily antes de aplicar
% el comando de estilo (#2). Sin declarar, \toc@font es un no-op.
% El folio (#5) hereda \toc@font porque va dentro del mismo \begingroup.

$if(titulos-sans)$
\newcommand{\toc@font}{\sffamily}
$else$
\newcommand{\toc@font}{}
$endif$

\newcommand{\toc@line}[5]{%
  % #1=leftindent  #2=fontcmd  #3=number  #4=title  #5=page
  \begingroup
    \toc@font
    \leftskip=#1\relax
    \parindent=\z@
    \rightskip=\z@
    \parfillskip=\z@ plus 1fil
    \leavevmode
    \if\relax\detokenize{#3}\relax\else
      \llap{#2{#3}\kern1em}%
    \fi
    #2{#4}%
    \nobreak\enspace\textperiodcentered\nobreak\enspace\normalfont#5%
    \par
  \endgroup
}

% ── Nivel 1 ───────────────────────────────────────────────────────────────────

\renewcommand*{\l@section}[2]{%
  \ifnum\c@tocdepth>\z@
    \addpenalty\@secpenalty
    \addvspace{\cftbeforesecskip}%
    \begingroup
      \def\numberline##1{}%
      \toc@line{0pt}{\MakeUppercase}{}{#1}{#2}%
    \endgroup
  \fi}

% ── Nivel 2 ───────────────────────────────────────────────────────────────────

\renewcommand*{\l@subsection}[2]{%
  \ifnum\c@tocdepth>\@ne
    \addvspace{\cftbeforesubsecskip}%
    \begingroup
      \def\numberline##1{}%
      \toc@line{0pt}{\upshape}{}{#1}{#2}%
    \endgroup
  \fi}

% ── Nivel 3 ───────────────────────────────────────────────────────────────────

\renewcommand*{\l@subsubsection}[2]{%
  \ifnum\c@tocdepth>\tw@
    \addvspace{\cftbeforesubsubsecskip}%
    \begingroup
      \def\numberline##1{}%
      \toc@line{\tocsubindent}{\itshape}{}{#1}{#2}%
    \endgroup
  \fi}

% ── Entradas sin folio (via \addcontentsline{toc}{part}{Título}) ──────────────

\providecommand{\l@part}[2]{}
\renewcommand{\l@part}[2]{%
  \addpenalty{-\@highpenalty}%
  \addvspace{0.5\baselineskip}%
  \noindent\normalfont #1\par}

\makeatother
